\section{Introduction}
Online experiments on social networks and marketplaces violate the no-interference assumption: a user's outcome depends on the treatments of the users she is connected to. The quantity a platform usually wants is the \emph{global average treatment effect} (GATE), the mean outcome difference between ``everyone treated'' and ``nobody treated''. Because this is a counterfactual that is never observed in a single experiment, every estimator must extrapolate. The difference-in-means (DM) estimator extrapolates not at all, and is biased by the spillover it ignores. Design-based estimators (Horvitz--Thompson/H\'ajek on exposure conditions \citep{aronow2017,ugander2013}) are unbiased for a coarsened target but have large variance when full-neighbourhood exposure is rare. A third family, \emph{regression-based} estimators \citep{chin2019,han2023}, fits a model from the assignment (and graph-derived exposure features) to the outcome, and then predicts the outcome under all-treated and all-control assignments. These are attractive because they are cheap, use all units, and need no exposure probabilities. But the prediction points ``all treated'' and ``all control'' usually lie far outside the observed feature distribution, and it is not well understood \emph{how accurate} such estimators are, \emph{what determines} their error, and \emph{when they fail}.

We study this question with a controlled simulation study backed by an exact error decomposition. Our contributions are:
\begin{enumerate}[leftmargin=*,itemsep=0pt,topsep=1pt]
\item \textbf{Exact decomposition and design diagnostics (\S\ref{sec:theory}).} For least-squares GATE estimators we show $\hat\tau=\sum_i w_iy_i$ with weights $w$ computable from the realised assignment alone. The error splits exactly into a noise term with variance $\sigma^2\|w\|_2^2$ and a misspecification/extrapolation term controlled by $\|w\|_1$ and the approximation residual. On regular graphs under Bernoulli($p$) designs, $n\|w\|_2^2\to(d+1)/(p(1-p))$, so error grows like $\sqrt{d/n}$, and the large-degree bias equals the \emph{local slope} $g'(p)$ of the exposure--response curve rather than its secant slope.
\item \textbf{A large, reproducible simulation study (\S\ref{sec:exp})} over eight interference mechanisms (linear, saturating, convex, threshold, two-hop, heterogeneous, equilibrium peer effects), six graphs (ER, BA, WS, SBM, and the public ego-Facebook and Facebook100-Stanford3 graphs), Bernoulli, cluster and randomised-saturation designs, with eleven estimators including a gradient-boosted learner and an H\'ajek baseline. All results report Monte-Carlo uncertainty.
\item \textbf{Findings, positive and negative.} The noise term is predicted almost exactly; the bias is \emph{not} reduced by more data; cluster randomisation helps only when the graph has good clusters; flexible learners are worse than linear regression under Bernoulli designs because of extrapolation; richer linear models overfit; neighbourhood misspecification and standard error estimators are examined, with failures reported (\S\ref{sec:res}).
\end{enumerate}
